\subsection{具最高权的单模}

回顾固定 B 在 \(h_{Q}^{*}\) 上定义了一个序关系（第六章 §1 第 6 节）。\(h_{Q}^{*}\) 中 \(\geq 0\) 的元素是 B 的元素以 \(\geq 0\) 的有理数为系数的线性组合。

更一般地，我们将考虑元素 \(\lambda, \mu \in \mathfrak{h}^*\) 之间的如下序关系：

\(\lambda - \mu\) 是 B 的元素以 \(\geq 0\) 的有理数为系数的线性组合。

引理 2. 设 V 是一个单 g-模，\(\omega\) 是 V 的一个权。下列条件等价：

\begin{enumerate}
\def\labelenumi{(\roman{enumi})}
\item
  V 的每个权都形如 \(\omega -\mu\) ，其中 \(\mu\) 是一个 \(\geq 0\) 的根权；
\item
  \(\omega\) 是 V 的最高权；
\item
  对一切 \(\alpha \in \mathrm{B}\) ，\(\omega + \alpha\) 不是 V 的权；
\item
  存在权为 \(\omega\) 的本原元。
\item
  \(\Longrightarrow\) (ii) \(\Longrightarrow\) (iii): 这是显然的。
\item
  \(\Longrightarrow\) (iv): 假设条件 (iii) 满足。对一切 \(h\in \mathfrak{h}\) ，
\end{enumerate}

\(\operatorname {Ker}(h_{\mathrm{V}} - \omega (h))\)

是非零的，包含于 \(V^{\omega}\) ，且在 \(h_{V}\) 下稳定。对 dim h 用归纳法，我们看到存在 \(V^{\omega}\) 中的非零 v 使 \(h v \subset kv\) 。条件 (iii) 蕴含 \(n_{+}v = 0\) ，故 v 是本原的。

\begin{enumerate}
\def\labelenumi{(\roman{enumi})}
\setcounter{enumi}{3}
\tightlist
\item
  \(\Longrightarrow\) (i): 设 \(v\) 是权为 \(\omega\) 的本原元。由于 V 是单的，V 由 \(v\) 生成（作为 \(\mathfrak{g}\) -模）。断言 (i) 于是由命题 1 得出。
\end{enumerate}

证毕。

于是，对任何单 g-模，本原元的存在性等价于最高权的存在性，或等价于极大权的存在性。

存在这样的单 \(\mathfrak{sl}(2,\mathbf{C})\) -模 V：对任何嘉当子代数 \(\mathfrak{h}\) （\(\mathfrak{sl}(2,\mathbf{C})\) 的）它都没有权（§1 习题 14 \(f\) )）。这些模在 \(\mathbf{C}\) 上是无限维的（§1 第 3 节 定理 1）。

命题 3. 设 V 是一个单 \(\mathfrak{g}\) -模，具有最高权 \(\omega\) 。

\begin{enumerate}
\def\labelenumi{(\roman{enumi})}
\item
  V 的本原元就是 \(V^{\omega}\) 的非零元素。
\item
  V 作为 \(\mathfrak{h}\) -模是半单的。
\item
  我们有 \(V = \bigoplus_{\lambda \in h^{*}} V^{\lambda}\) 。对一切 \(\lambda \in h^{*}\) ，\(V^{\lambda}\) 是有限维的。我们有 \(\dim V^{\omega} = 1\) 。
\item
  \(\mathfrak{g}\) -模 V 是绝对单的。
\end{enumerate}

断言 (i)、(ii) 与 (iii) 由命题 1 与引理 2 得出。断言 (iv) 由命题 1 (iii) 与代数第八章 §7 第 3 节得出。

推论. 若 \(\mathrm{V}\) 是有限维的，则典则同态 \(\mathrm{U}(\mathfrak{g}) \to \operatorname{End}(\mathrm{V})\) 是满射。

这由 (iv) 得出，参见代数第八章 §3 第 3 节。

命题 4. 设 V 是一个具有最高权 \(\omega\) 的单 g-模，X 是一个半单 g-模，\(X'\) 是 X 中 V 型的同构型分支。则 \(X'\) 是 X 的由 \(X_{\pi}^{\omega}\) 生成的子模。它的长度等于 \(X_{\pi}^{\omega}\) 的维数。

设 \(X''\) 是 X 的由 \(X_{\pi}^{\omega}\) 生成的子模。显然 X 的每个同构于 V 的子模都包含在 \(X''\) 中。故 \(X' \subset X''\) 。另一方面，设 \(\varPhi\) 是从 g-模 V 到 g-模 X 的同态组成的集合。\(X'\) 的长度是 \(\dim_k \varPhi\) （代数第八章 §4 第 4 节），即 \(\dim_k X_{\pi}^{\omega}\) （命题 2）。

